% Cyclic-inhibition figure, Figure S1.1 of the Supplementary Material.
%
% Reads ../results/ring-timeseries.dat, ../results/ring-bifurcation.dat and
% ../results/ring-threshold.tex, all written by ring_cycle.py.  Build from the
% repository root with
%
%     python3 ring_cycle.py --replot
%
% which runs latexmk in this directory and moves the PDF to ../figures/.
\documentclass[border=2pt]{standalone}
\usepackage{amsmath}
\usepackage{pgfplots}
\pgfplotsset{compat=1.18}
\usepgfplotslibrary{groupplots}

\def\ringcstar{2.0547}%


\definecolor{popone}{HTML}{1F4E79}
\definecolor{poptwo}{HTML}{C1121F}
\definecolor{popthree}{HTML}{2A9D8F}

\pgfplotsset{
    ringaxis/.style={
        width=7.4cm, height=5.2cm,
        scale only axis,
        axis lines=left,
        line width=0.5pt,
        tick align=outside,
        tick label style={font=\footnotesize},
        label style={font=\small},
        title style={font=\small, yshift=-2pt},
        legend style={font=\footnotesize, draw=none, fill=none,
                      inner sep=1pt, row sep=-1pt},
        every axis plot/.append style={line width=0.9pt},
    },
}

\begin{document}
\begin{tikzpicture}
\begin{groupplot}[
    group style={group size=2 by 1, horizontal sep=1.9cm},
    ringaxis,
]

% ---- (a) the attracting periodic orbit at c = 5 ---------------------
\nextgroupplot[
    xlabel={$t$}, ylabel={activity},
    title={(a) periodic orbit at $c=5$},
    xmin=0, xmax=12, ymin=-0.75, ymax=0.95,
    ytick={-0.6,-0.3,0,0.3,0.6},
    legend pos=north east, legend columns=3,
]
\addplot[popone]   table[x=t, y=x1] {../results/ring-timeseries.dat};
\addplot[poptwo]   table[x=t, y=x2] {../results/ring-timeseries.dat};
\addplot[popthree] table[x=t, y=x3] {../results/ring-timeseries.dat};
\legend{$x_1$, $x_2$, $x_3$}

% ---- (b) bifurcation diagram ----------------------------------------
\nextgroupplot[
    xlabel={$c$}, ylabel={$x_1$ on the attractor},
    title={(b) onset of oscillation},
    xmin=1, xmax=10, ymin=-0.75, ymax=0.95,
    ytick={-0.6,-0.3,0,0.3,0.6},
    legend style={at={(0.98,0.30)}, anchor=east}, legend columns=1,
]
\addplot[black, dashed, line width=0.6pt, forget plot]
    coordinates {(\ringcstar,-0.75) (\ringcstar,0.95)};
% the attractor first, the equilibrium branch on top of it, so that the two
% remain distinguishable where they coincide, that is below c*
\addplot[popone] table[x=c, y=xmax] {../results/ring-bifurcation.dat};
\addlegendentry{attractor}
\addplot[popone, forget plot] table[x=c, y=xmin] {../results/ring-bifurcation.dat};
\addplot[gray, dash pattern=on 3pt off 2pt]
    table[x=c, y=xeq] {../results/ring-bifurcation.dat};
\addlegendentry{equilibrium $x^\ast$}
\node[anchor=west, font=\footnotesize] at (axis cs:\ringcstar,0.85)
    {\,$c^\ast=\ringcstar$};
\end{groupplot}
\end{tikzpicture}
\end{document}
